\noindent HENRY, Marco Antônio Cottorello. NEVES, Murilo de Souza. SILVA JUNIOR, Rodrigo Macedo da.
\textbf{Virtualization, Cloud and Containers}. 2026. 24 pages.
(Graduação em Ciência da Computação) -- State University of Londrina,
Londrina, 2026.

\vspace{2\baselineskip}
%Abstract
\begin{center}
  \textbf{ABSTRACT}
\end{center}

\vspace{\baselineskip}

\noindent The consolidation of cloud computing has shifted information technology infrastructure from dedicated servers to resources provisioned on demand and supported by successive layers of abstraction. This work presents that stack in an integrated manner, starting from the canonical NIST definition, with its five essential characteristics, three service models and four deployment models, and from the criterion that actually separates IaaS, PaaS and SaaS, namely the position of the control boundary over the software stack. It then examines two implementations located in the resource abstraction and control layer. The first is OpenStack, whose modular architecture and whose identity, image, compute and networking services turn physical servers into infrastructure provisionable through a programmable interface. The second is Docker, whose decomposed engine, kernel isolation primitives and local network models define the container as an isolated process running on a shared kernel. On that basis, the work analyses the attack surface arising from kernel sharing, covering privilege containment mechanisms, image vulnerabilities and protection of inter-container traffic, and discusses energy efficiency in data centres, covering the PUE metric, dynamic server consolidation and thermal management. It concludes that isolation choices and consolidation density are interdependent, so that security and energy efficiency are coupled dimensions of a single infrastructure design.

\vfill

\noindent\textbf{Key words:} Cloud Computing. Virtualization. Containers. OpenStack. Docker. Energy Efficiency.
